\subsection{\texorpdfstring{\(p\)-free subsets and \(p\)-bases}{p-free subsets and p-bases}}

DEFINITION 1. --- Let K be a field and L a p-radical extension of height ≤1 of K. A family \((x_i)_{i \in I}\) of elements of L is said to be p-free over K (resp. a p-basis of L over K) if \(x_i \notin K\) for all \(i \in I\) and the homomorphism of \(\otimes_{i \in I} K(x_i)\) into L

derived from the canonical injections \(u_{i}:\mathbf{K}(x_{i})\to \mathbf{L}\) is injective (resp. bijective).

If \(a\) is an element of \(\mathbf{L} - \mathbf{K}\) , it is \(p\) -radical of height 1, and so its minimal polynomial over \(\mathbf{K}\) is \(\mathbf{X}^p - a^p\) (V, p.~24, Prop. 1); therefore \(\{1, a, \ldots, a^{p-1}\}\) is a basis of \(\mathbf{K}(a)\) over \(\mathbf{K}\) . Let \((x_i)_{i \in I}\) be a family of elements of \(\mathbf{L} - \mathbf{K}\) and \(\Lambda\) the subset of \(\mathbf{N}^{(I)}\) consisting of all families of finite support \(\alpha = (\alpha_i)_{i \in I}\) such that \(\alpha_i < p\) for all \(i \in I\) ; Prop. 9 (III, p.~471) shows that the elements \(\otimes_{i \in I} x_i^{\alpha_i}\) as \(\alpha\) runs

over \(\Lambda\) form a basis of the vector space \(\otimes_{i\in I}\mathbf{K}(x_i)\) over \(\mathbf{K}\) . Moreover, the canonical

homomorphism of \(\otimes_{i\in I}\mathbf{K}(x_i)\) into \(L\) has as image \(\mathbf{K}(x_i)_{i\in I}\) (V, p.~18, Cor. 1), we

thus have the following proposition :

PROPOSITION 1. --- Let K be a field, L a p-radical extension of height ≤1 of K and \(\mathbf{x} = (x_{i})_{i \in I}\) a family of elements of L. Then the vector space \(\mathrm{K}(x_{i})_{i \in I}\) over K is generated by the products \(x^{\alpha} = \prod_{i \in I} x_{i}^{\alpha_{i}}\) for \(\alpha\) in \(\Lambda\) . For

\((x_{i})_{i\in I}\) to be \(p\) -free (resp. a \(p\) -basis) it is necessary and sufficient that the family \((\mathbf{x}^{\alpha})_{\alpha \in \Lambda}\) should be free over \(\mathbf{K}\) (resp. a basis of \(\mathbf{L}\) over \(\mathbf{K}\) ).

COROLLARY. --- Let L' be an extension of a field K, generated by two linearly disjoint subextensions K' and L. Suppose that L is p-radical of height \(\leqslant 1\) over K; then L' is p-radical of height \(\leqslant 1\) over K'. Moreover, for a family of elements of L to be p-free over K (resp. a p-basis of L over K) it is necessary and sufficient that it should be p-free over K' (resp. a p-basis of L' over K').

We have \(L^{p} \subset K \subset K'\) and \(L' = K'(L)\) , whence \(L'^{p} = K'^{p}(L^{p}) \subset K'\) . In other words, \(L'\) is a p-radical extension of height \(\leqslant 1\) of \(K'\) . The other assertions of the corollary follow at once from Prop. 1 and from V, p.~14, Prop. 5.

Remark. --- Let K be a field, L a p-radical extension of height \(\leqslant1\) of K and \((x_{i})_{i\in I}\) a family of elements of L. Let A be the polynomial algebra \(\mathrm{K}[(\mathrm{X}_{i})_{i\in I}]\) , a the ideal of A generated by the polynomials \(X_{i}^{p}-x_{i}^{p}\) and \(\varphi:A\to L\) the K-algebra homomorphism such that \(\varphi(\mathrm{X}_{i})=x_{i}\) for each \(i\in I\) . The family \((x_{i})_{i\in I}\) is p-free if and only if the kernel of \(\varphi\) is equal to a. For \(\mathrm{K}[(\mathrm{X}_{i})]/\mathfrak{a}\) may be identified with the algebra \(\otimes\mathrm{K}[\mathrm{X}_{i}]/(\mathrm{X}_{i}^{p}-x_{i}^{p})\) .

Let L be a p-radical extension of height \(\leqslant1\) of the field K. A subset S of L is said to be p-free (resp. a p-basis) if the family defined by the identity mapping of S onto itself is p-free (resp. a p-basis). For a family \((x_{i})_{i\in I}\) of elements of L to be p-free (resp. a p-basis) it is necessary and sufficient that the mapping \(i\mapsto x_{i}\) should be a bijection of I on a p-free subset (resp. a p-basis) of L. By Prop. 1, every subset of a p-free set is p-free; conversely, if S is a subset of L all of whose finite subsets are p-free, then S is p-free. Finally, a p-basis of L over K is a p-free subset B such that \(\mathrm{L}=\mathrm{K}(\mathrm{B})\) .

PROPOSITION 2. --- Let K be a field, L a p-radical extension of height ≤1 of K and S a subset of L. For S to be p-free it is necessary and sufficient that K(T) ≠ K(S) for every subset T ≠ S of S.

Suppose first that S is p-free and let \(T \neq S\) be a subset of S. Denote by \(\Lambda\) the set of families with finite support \(\alpha = (\alpha_{s})_{s \in S}\) of integers between 0 and p - 1; let \(\Lambda'\) be the subset of \(\Lambda\) consisting of the families \(\alpha = (\alpha_{s})_{s \in S}\) such that \(\alpha_{s} = 0\) for all s in S - T. We also put \(u_{\alpha} = \prod_{s \in S} s^{\alpha_{s}}\) for \(\alpha \in \Lambda\) . Then (V, p.~98,

Prop. 1) the family \((u_{\alpha})_{\alpha \in \Lambda}\) is a basis of \(\mathbf{K}(\mathbf{S})\) over \(\mathbf{K}\) and the subfamily \((u_{\alpha})_{\alpha \in \Lambda'}\) is a basis of \(\mathbf{K}(\mathbf{T})\) over \(\mathbf{K}\) . Since \(\Lambda' \neq \Lambda\) , we have \(\mathbf{K}(\mathbf{T}) \neq \mathbf{K}(\mathbf{S})\) .

Suppose now that S is not \(p\) -free. Then there exists an integer \(n \geqslant 1\) and a sequence of elements \(x_1, \ldots, x_n\) of S such that \((x_1, \ldots, x_{n-1})\) is \(p\) -free but \((x_1, \ldots, x_n)\) is not. We have \([\mathbf{K}(x_1, \ldots, x_{n-1}) : \mathbf{K}] = p^{n-1}\) and \([\mathbf{K}(x_1, \ldots, x_n) : \mathbf{K}] < p^n\) . Since \([\mathbf{K}(x_1, \ldots, x_n) : \mathbf{K}]\) is a multiple of \([\mathbf{K}(x_1, \ldots, x_{n-1}) : \mathbf{K}]\) , we thus have

\[
[ \mathrm{K} (x _ {1}, \dots , x _ {n}): \mathrm{K} ] = [ \mathrm{K} (x _ {1}, \dots , x _ {n - 1}): \mathrm{K} ].
\]

whence \(x_{n} \in \mathbf{K}(x_{1}, \ldots, x_{n-1})\) . This shows that \(\mathbf{K}(\mathbf{S} - \{x_{n}\}) = \mathbf{K}(\mathbf{S})\) .

PROPOSITION 3. --- Let K be a field, L a p-radical extension of height ≤1 of K and S, T two subsets of L. The following conditions are equivalent :

\begin{enumerate}
\def\labelenumi{\alph{enumi})}
\item
  The subset \(\mathbf{S}\) is \(p\) -free over \(\mathbf{K}\) and \(\mathbf{T}\) is \(p\) -free over \(\mathbf{K}(\mathbf{S})\) .
\item
  \(\mathbf{S} \cap \mathbf{T} = \varnothing\) and \(\mathbf{S} \cup \mathbf{T}\) is \(p\) -free over \(\mathbf{K}\) .
\end{enumerate}

If T is p-free over \(\mathbf{K}(\mathbf{S})\) , we have \(\mathbf{T} \cap \mathbf{K}(\mathbf{S}) = \varnothing\) and a fortiori \(S \cap T = \varnothing\) . We may therefore assume S and T disjoint. Let \(\Lambda\) be the subset of \(\mathbf{N}^{(S \cup T)}\) consisting of all families \(\alpha = (\alpha_x)_{x \in S \cup T}\) with \(\alpha_x < p\) for all \(x \in S \cup T\) , and define subsets \(\Lambda'\) of \(\mathbf{N}^{(S)}\) and \(\Lambda''\) of \(\mathbf{N}^{(T)}\) in similar fashion. We can in a natural way, identify \(\mathbf{N}^{(S \cup T)}\) with \(\mathbf{N}^{(S)} \times \mathbf{N}^{(T)}\) , and then \(\Lambda\) is identified with \(\Lambda' \times \Lambda''\) . For \(\alpha \in \Lambda\) put \(u_\alpha = \prod_{x \in S \cup T} x^{\alpha_x}\) , and similarly define \(u_\beta'\) and \(u_\gamma''\) for \(\beta \in \Lambda'\) and

\(\gamma \in \Lambda''\) . We have \(u_{\alpha} = u_{\beta}' u_{\gamma}''\) for \(\alpha = (\beta, \gamma)\) in \(\Lambda = \Lambda' \times \Lambda''\) . Moreover \((u_{\beta}')_{\beta \in \Lambda'}\) generates the vector K-space K(S) (V, p.~94, Prop. 1). For S \(\cup\) T to be \(p\) -free over K it is necessary and sufficient that the family \((u_{\beta}' u_{\gamma}'')_{\beta \in \Lambda', \gamma \in \Lambda''}\) should be free over K (V, p.~98, Prop. 1); it comes to the same (II, p.~222) to suppose that the family \((u_{\beta}')_{\beta \in \Lambda'}\) is free over K and the family \((u_{\gamma}'')_{\gamma \in \Lambda''}\) free over K(S). Now the equivalence of \(a\) ) and \(b\) ) follows from Prop. 1 (V, p.~98).

COROLLARY. --- Let L be a p-radical extension of height ≤1 of K and M a subextension of L. Then L is p-radical of height ≤1 over M and M is p-radical of height ≤1 over K. Moreover, if B is a p-basis of M over K and C is a p-basis of L over M, then B ∩ C = ∅ and B ∪ C is a p-basis of L over K.

Let K be a field. A family \((x_{i})_{i\in I}\) is said to be a (absolute) p-basis of K if it is a p-basis of K over \(K^{p}\) . For every integer \(n\geqslant1\) we denote by \(\Lambda(n)\) the subset of \(\mathbf{N}^{(I)}\) consisting of all \(\alpha=(\alpha_{i})_{i\in I}\) such that \(\alpha_{i}<p^{n}\) for all \(i\in I\) .

PROPOSITION 4. --- Let \(\mathbf{x} = (x_i)_{i \in I}\) be a \(p\) -basis of K. For every integer \(n \geqslant 1\) , the family \((\mathbf{x}^\alpha)_{\alpha \in \Lambda(n)}\) is a basis of K over \(\mathbf{K}^{p^n}\) .

For n = 1 the assertion reduces to Prop. 1 (V, p.~98). The set \(\Lambda(n)\) consists of elements of \(\mathbf{N}^{(I)}\) of the form \(\alpha = \beta + p^{n-1}\gamma\) with \(\beta \in \Lambda(n-1)\) and \(\gamma \in \Lambda(1)\) . Such a decomposition is unique and we have \(\mathbf{x}^{\alpha} = \mathbf{x}^{\beta}(\mathbf{x}^{\gamma})^{p^{n-1}}\) . Moreover, the family \((x_{i}^{p^{n-1}})_{i \in I}\) is clearly a p-basis of \(K^{p^{n-1}}\) over \(K^{p^n}\) , hence the family \((\mathbf{x}^{\gamma})^{p^{n-1}}\) is a basis of \(K^{p^{n-1}}\) over \(K^{p^n}\) . We thus obtain the conclusion by induction, using II, p.~222, Prop. 25.

